% reports/ieee_report/sections/p2_architecture.tex
% Author: Vani Rudra (Task P2_03 — Autoencoder Architecture & Theory)
% IEEE Conference Paper Section — Card Fraud Triage (Phase 2)

\section{Deep Autoencoder Architecture \& Theory}
\label{sec:p2_architecture}

To model the non-linear manifold of legitimate credit-card transactions without
relying on delayed fraud labels, we implement both a symmetric Undercomplete
Autoencoder (AE) and a Denoising Autoencoder (DAE).  Consistent with Phase~1
(Sec.~\ref{sec:p2_dl_data_pipeline}), all neural models are trained
exclusively on the $170{,}524$ legitimate training transactions
($\mathbf{X}_{\text{train,legit}} \in \mathbb{R}^{170524 \times 31}$),
early-stopped strictly on the $56{,}904$ validation-legitimate records, and
scored by per-sample Mean Squared Error (MSE) so that higher values denote
greater abnormality.

% ─── 1. Formal Formulation ───────────────────────────────────────────────────
\subsection{Encoder--Decoder Formulation and Objective}
\label{sec:ae_formulation}

Let $\mathbf{x} \in \mathbb{R}^{d}$ ($d = 31$ for the base feature set) denote
a standardized transaction vector.  An autoencoder consists of an encoder
$f_{\boldsymbol{\theta}} : \mathbb{R}^{d} \to \mathbb{R}^{k}$ mapping input
vectors into a low-dimensional bottleneck code
$\mathbf{z} = f_{\boldsymbol{\theta}}(\mathbf{x})$, and a symmetric decoder
$g_{\boldsymbol{\phi}} : \mathbb{R}^{k} \to \mathbb{R}^{d}$ producing the
reconstruction $\hat{\mathbf{x}} = g_{\boldsymbol{\phi}}(\mathbf{z})$.
Parameters $(\boldsymbol{\theta}, \boldsymbol{\phi})$ are learned by minimizing
the expected reconstruction MSE over the legitimate training distribution:
\begin{equation}
    \min_{\boldsymbol{\theta}, \boldsymbol{\phi}} \;
    \mathbb{E}_{\mathbf{x} \sim P_{\text{legit}}}
    \left[ \frac{1}{d} \, \bigl\| \mathbf{x} - g_{\boldsymbol{\phi}}\bigl(f_{\boldsymbol{\theta}}(\mathbf{x})\bigr) \bigr\|_2^2 \right]
    + \lambda \, \|\boldsymbol{\theta}\|_2^2 + \lambda \, \|\boldsymbol{\phi}\|_2^2,
    \label{eq:ae_objective}
\end{equation}
where $\lambda$ is the AdamW/Adam weight-decay coefficient.  The encoder uses
hidden layer widths $[h_1, h_2]$ (default $[24, 16]$) with ReLU-family
activations initialized via He/Kaiming normal initialization~\cite{he2015delving},
followed by a bottleneck of width $k < d$.  The decoder mirrors the encoder
($k \to h_2 \to h_1 \to d$) and terminates in an unconstrained \emph{linear}
output layer because standardized input features are real-valued and unbounded.

At inference time, each transaction receives an anomaly score equal to its
per-sample MSE across coordinates:
\begin{equation}
    s_{\text{AE}}(\mathbf{x}) = \frac{1}{d} \sum_{j=1}^{d} \bigl(x_j - \hat{x}_j\bigr)^2,
    \label{eq:ae_score}
\end{equation}
while individual coordinate residuals $e_j(\mathbf{x}) = (x_j - \hat{x}_j)^2$
provide feature-level interpretability for flagged alerts.  Additionally, the
bottleneck representation $\mathbf{z} = f_{\boldsymbol{\theta}}(\mathbf{x}) \in
\mathbb{R}^{k}$ is exposed via the \texttt{encode()} interface for downstream
latent-space GMM density estimation ($H_1$) and hybrid scoring ($H_2, H_3$).

% ─── 2. Undercompleteness, Bias-Variance, and PCA Equivalence ────────────────
\subsection{Undercompleteness, PCA Equivalence, and Nonlinearity}
\label{sec:ae_pca_theory}

\textbf{Why the bottleneck must be undercomplete ($k < d$).}
If $k \ge d$ (or if network capacity is unconstrained), minimizing
Eq.~\eqref{eq:ae_objective} permits the trivial identity mapping
$g_{\boldsymbol{\phi}}(f_{\boldsymbol{\theta}}(\mathbf{x})) \approx \mathbf{x}$
across the entire ambient space $\mathbb{R}^d$.  In that degenerate regime,
both legitimate transactions and fraudulent outliers are reconstructed with
near-zero error, collapsing anomaly contrast.  Enforcing $k < d$ creates an
information bottleneck: the network has only $k$ degrees of freedom to encode
normal spending correlations, forcing high error on off-manifold transactions.

\textbf{Bias--variance trade-off in $k$.}
Selecting $k$ balances reconstruction fidelity against anomaly discrimination:
\begin{itemize}
    \item \emph{Too small ($k = 2$):} High bias / underfitting — legitimate
          transactions cannot be faithfully reconstructed, elevating baseline
          noise and obscuring genuine fraud spikes.
    \item \emph{Too large ($k \ge 16$):} Low bias / over-generalization — the
          decoder captures high-order variation including regions populated by
          fraudulent records, reducing separation.
\end{itemize}
We sweep $k \in \{2, 4, 8, 12, 16\}$ on the validation set to locate the
optimal operating point.

\textbf{Linear AE vs.\ PCA and the role of nonlinearity.}
When all activations in $f_{\boldsymbol{\theta}}$ and $g_{\boldsymbol{\phi}}$
are linear and trained under MSE, Baldi and Hornik~\cite{baldi1989neural}
showed that the global minimum of Eq.~\eqref{eq:ae_objective} has a unique
subspace solution: the encoder/decoder weights span the exact same $k$-dimensional
principal subspace as Principal Component Analysis (PCA).  We verify this
analytically in Sec.~\ref{sec:p2_ae_results} as a sanity check against our
Phase~1 PCA baseline.  Replacing linear projections with ReLU, LeakyReLU, or
Tanh activations enables the autoencoder to learn curved, piecewise-linear
(or smooth) non-linear manifolds~\cite{hinton2006reducing} rather than a single
flat hyperplane.

% ─── 3. Denoising Autoencoder & Score Matching ───────────────────────────────
\subsection{Denoising Autoencoders and Manifold Score Matching}
\label{sec:dae_theory}

Even with $k < d$, a deep multi-layer autoencoder can memorize local neighborhoods
around training points.  A Denoising Autoencoder (DAE)~\cite{vincent2008extracting}
explicitly regularizes the representation by corrupting each clean training
input with isotropic Gaussian noise during training:
\begin{equation}
    \tilde{\mathbf{x}} = \mathbf{x} + \sigma \boldsymbol{\varepsilon},
    \qquad \boldsymbol{\varepsilon} \sim \mathcal{N}(\mathbf{0}, \mathbf{I}_d),
    \quad \sigma \in \{0.05, 0.1, 0.2, 0.5\},
\end{equation}
and minimizing the reconstruction error against the \emph{clean} target:
\begin{equation}
    \mathcal{L}_{\text{DAE}}(\boldsymbol{\theta}, \boldsymbol{\phi})
    = \mathbb{E}_{\mathbf{x}, \boldsymbol{\varepsilon}}
    \left[ \frac{1}{d} \, \bigl\| \mathbf{x} - g_{\boldsymbol{\phi}}\bigl(f_{\boldsymbol{\theta}}(\tilde{\mathbf{x}})\bigr) \bigr\|_2^2 \right].
    \label{eq:dae_loss}
\end{equation}
At inference time, no noise is injected ($\sigma = 0$), and the score remains
Eq.~\eqref{eq:ae_score}.  As proven by Vincent~\cite{vincent2011connection},
training a DAE with small Gaussian noise $\sigma$ approximates score matching:
the residual vector $g_{\boldsymbol{\phi}}(f_{\boldsymbol{\theta}}(\mathbf{x})) - \mathbf{x}$
estimates $\sigma^2 \nabla_{\mathbf{x}} \log p(\mathbf{x})$, pulling corrupted
points back onto the high-density normal manifold.

% ─── 4. When Reconstruction Error Fails ──────────────────────────────────────
\subsection{Failure Modes of Reconstruction-Error Scoring}
\label{sec:ae_failure_modes}

While MSE reconstruction error effectively flags \emph{off-manifold} anomalies,
it exhibits two structural blind spots that motivate Phase~2 hybrid models:
\begin{enumerate}
    \item \textbf{In-manifold (on-subspace) anomalies:} A fraudulent transaction
          whose feature vector happens to lie on or near the learned $k$-dimensional
          manifold — for example, a legitimate-looking coordinate combination in a
          sparse tail region — achieves small $\|\mathbf{x} - \hat{\mathbf{x}}\|_2^2$
          and goes undetected.  Modeling latent density $-\log p(\mathbf{z})$
          directly addresses this gap.
    \item \textbf{High-variance coordinate dominance:} Unweighted MSE scales
          proportionally to coordinate variance $\operatorname{Var}(X_j)$.
          Without train-legit standardization, heavy-tailed features (such as raw
          \texttt{Amount}) dominate both gradients and anomaly scores, masking
          discriminative deviations in lower-variance PCA components.
\end{enumerate}

% ─── 5. Optimization & Regularization Protocol ───────────────────────────────
\subsection{Training and Regularization Protocol}
\label{sec:ae_training_protocol}

All models are trained with Adam~\cite{kingma2015adam} (initial learning rate
$10^{-3}$, batch size $512$, max $100$ epochs), gradient-norm clipping at $5.0$,
and \texttt{ReduceLROnPlateau} (factor $0.5$, patience $5$).  Early stopping
monitors \emph{validation-legitimate} MSE exclusively (patience $10$), ensuring
fraud labels never influence gradient updates or checkpoint stopping.  We
systematically ablate $L_2$ weight decay, dropout~\cite{srivastava2014dropout},
Gaussian input noise (DAE), batch normalization, early stopping, and $L_1$
latent sparsity to evaluate generalization and train--validation gap.

